%! Author = zhiweizhang
%! Date = 10.06.25

\chapter{Discussion and Analysis}\label{chapter:discussion-and-analysis}

This chapter reflects critically on the performance of the evaluated models, connecting the empirical results to the
research objectives outlined in ~\autofullref{chapter:introduction}. The primary aim of this thesis was to develop
and compare image classification models capable of accurately recognizing and categorizing Nazi symbols from varied
visual inputs. Through detailed analysis of both binary and multi-class classification results, this chapter
interprets key trends, identifies strengths and limitations of different model families, and discusses the practical
and ethical implications of deploying such systems.

\section{Analysis of Binary Classification Performance}

The results presented in ~\autoref{tab:binary-results} reveal a wide range of performances across different models.
Notably, classical \ac{ML} classifiers leveraging OpenCLIP embeddings consistently ranked among the top
performers. The best overall F1-score (0.939) was achieved by the \textbf{OpenCLIP + ~\ac{SVC}} model, which also
demonstrated high precision (0.948) and recall (0.929). Similarly, the \textbf{OpenCLIP + Logistic Regression}
model achieved the highest ~\ac{AUC} score (1.000) and highest \ac{AP} (0.964), along with a strong recall of 0.882
and an F1-score of 0.902, indicating a solid ability to correctly classify Nazi symbols while minimizing false
negatives.

Among \ac{DL} models, both \textbf{\ac{ViT}} and \textbf{\ac{YOLO}} performed competitively. \ac{ViT} achieved an
F1-score of 0.901, the highest among all standalone neural models, while \ac{YOLO} achieved an F1-score of 0.857 with
a notably high recall of 0.932, highlighting its capacity to correctly classify positive symbol instances.
\ac{ResNet} also performed well, with an F1-score of 0.867, slightly outperforming EfficientNet (0.767).

Anomalously, the \textbf{Moondream} model achieved an exceptionally high recall of 0.995 but suffered from very low
precision (0.153), resulting in a low F1-score of just 0.265. This suggests a strong tendency to overpredict the
Nazi class, leading to a high number of false positives. Similarly, the \textbf{OpenCLIP zero-shot} model performed
poorly, with a low F1-score of 0.025, highlighting the limitations of zero-shot classification without additional
fine-tuning or classifier integration.

Several classical ensemble models also showed robust performance when paired with OpenCLIP embeddings. For instance,
\textbf{OpenCLIP + VotingClassifier} reached an F1-score of 0.908 and precision of 0.945, while
\textbf{OpenCLIP + AdaBoostClassifier} and \textbf{OpenCLIP + GradientBoostingClassifier} achieved F1-scores of
0.842 and 0.838, respectively. However, the \textbf{OpenCLIP + RandomForestClassifier} produced an unusually high
precision (0.983) but very low recall (0.097), resulting in a poor F1-score (0.177), indicating that it was highly
conservative in predicting the Nazi class.

Overall, models that combine OpenCLIP embeddings with classical classifiers—especially \ac{SVC}, KNeighborsClassifier, and
VotingClassifier—outperformed all other approaches in this binary classification task. Nonetheless, end-to-end deep
learning models like \ac{ViT} and \ac{YOLO} remain competitive and demonstrate strong generalization capabilities,
particularly when trained on high-resolution inputs and sufficient labeled data.

\subsection{Binary Classification Confusion Matrix Analysis}

~\autoref{fig:confusion-matrix-binary} presents the confusion matrices for the binary classification task, where the
goal is to distinguish between \textit{Nazi} and \textit{Non-Nazi} images.

The OpenCLIP encoder combined with \ac{SVC} achieves the best overall balance, producing the
highest number of true positives (194) and true negatives (14,802), alongside the lowest false positives (11) and
false negatives (11). This indicates a strong capability to correctly classify Nazi symbols while minimizing
misclassifications.

\ac{YOLO} shows excellent sensitivity, with the second lowest false negative count (14), meaning it rarely misses Nazi
instances, but it has a relatively higher false positive rate (50), which might lead to more false alarms.

\ac{ResNet} also performs well, showing a favorable trade-off between false positives (34) and false negatives (22),
maintaining a good level of accuracy and reliability.

\ac{ViT} provides balanced performance, with moderate false positives (28) and false negatives (18), showing
robustness in both detecting Nazi symbols and correctly rejecting non-Nazi images.

EfficientNet, while demonstrating strong true negative results, has the highest false negatives (49), indicating
that it misses more Nazi cases compared to the other models, which might limit its applicability when recall is critical.

Overall, these confusion matrices confirm that the OpenCLIP + \ac{SVC} approach is the most effective for binary
classification in this context, offering a strong balance between precision and recall essential for sensitive
symbol classification.

\subsection{Binary Classification \ac{ROC} and \ac{PR} Curve Analysis}

The \ac{ROC} curve results indicate that most models achieved excellent discriminative ability in the binary
classification task, with several models reaching near-perfect \ac{AUC} values close to 1.0. Notably, OpenCLIP combined
with \ac{LR}, KNeighborsClassifier, and Random Forest classifiers, along with \ac{ViT}, \ac{YOLO},
and \ac{ResNet}, all demonstrated near-identical, near-optimal \ac{ROC} performance (\ac{AUC} $\approx$ 1.000),
suggesting their strong capacity to separate Nazi and Non-Nazi classes.

Ensemble methods such as AdaBoost and Gradient Boosting classifiers paired with OpenCLIP embeddings also performed
robustly (\ac{AUC} $\approx$ 0.998), while EfficientNet’s slightly lower \ac{AUC} of 0.993 still represents a high-quality
classification boundary. Bagging and Decision Tree classifiers showed comparatively lower \ac{AUC}s (0.955 and 0.933),
indicating reduced robustness and potential susceptibility to misclassification in borderline cases.

The \ac{AP} values from the \ac{PR} curves further underscore these findings, revealing that
OpenCLIP + \ac{LR} and \ac{ViT} provide the best balance between precision and recall, critical in scenarios where false
positives and false negatives carry significant consequences. The lower \ac{AP} of the OpenCLIP + Decision Tree
classifier (0.625) suggests that, despite reasonable \ac{AUC} scores, this model struggles to maintain precision at
higher recall levels, reflecting less reliable positive class predictions.

These combined results imply that embedding-based methods (OpenCLIP) paired with strong classifiers (\ac{LR},
KNeighborsClassifier) or transformer architectures (\ac{ViT}) are most effective for the task. In contrast, simpler or single-model
approaches such as Decision Trees show limitations in practical use despite acceptable \ac{ROC} metrics.

Only models that support probabilistic outputs are included in this analysis. Classifiers such as \texttt{\ac{SVC}},
\texttt{SGDClassifier}, and the \texttt{VotingClassifier}, as well as models like \texttt{Moondream} and
\texttt{OpenCLIP}, are omitted from \ac{ROC} and \ac{PR} curve evaluations, as discussed in
~\autoref{sec:evaluation-metrics}.

\subsection{Binary Classification with Reduced Nazi Symbol Classes}

To further examine model robustness and adaptability, a binary classification experiment was conducted. This
experiment was restricted to a reduced subset of Nazi symbols (\textit{siegrune}, \textit{swastika}, and
\textit{\ac{SS} skull}) against the full set of non-Nazi images. The goal was to evaluate model performance when the
semantic diversity of the positive class is reduced.

Despite the narrower scope, overall trends remained consistent with the full binary task (as showed in
~\autoref{tab:binary-reduced-symbols-result}). Transformer-based
architectures continued to perform strongly. \ac{ViT} achieved excellent results with an F1-score of \textbf{0.948},
indicating its strong ability to generalize even with reduced symbol variety. Among OpenCLIP-based hybrid models,
\textit{OpenCLIP + \ac{SVC}} achieved the best overall performance with an accuracy of \textbf{0.998},
precision of \textbf{0.973}, recall of \textbf{0.900}, and F1-score of \textbf{0.935}, while
\textit{OpenCLIP + KNeighborsClassifier} attained the highest recall at \textbf{0.989}.

Interestingly, \ac{YOLO} underperformed in this setup (F1-score \textbf{0.661}) compared to its results on the full
binary task, potentially due to its reliance on spatial object features that may be less effective when class
variety is reduced. Conversely, models like \textbf{Moondream} and \textbf{OpenCLIP} hybrids maintained high performance, indicating
that text-image aligned embeddings may still effectively capture distinguishing features in constrained settings.

As this is a focused experiment with a reduced symbol set, the evaluation was limited to the most informative and widely
supported metrics—accuracy, precision, recall, and F1-score. Metrics such as \ac{AUC} and \ac{AP} are omitted for
simplicity and consistency, especially given that not all models provide calibrated probability outputs.

This experiment highlights that reducing the complexity of the positive class does not necessarily simplify the task
uniformly for all models. While it benefits embedding-based methods and transformers, object detection models like
\ac{YOLO} may be affected due to loss of visual diversity.

\section{Analysis of Multi-Class Classification Performance}

~\autoref{tab:multi-class-results} summarizes the performance of all evaluated models on the multi-class
classification task. The \textbf{\ac{YOLO}} model achieved the best performance across all macro-averaged metrics,
with an accuracy of 0.884, precision of 0.865, recall of 0.818, and F1-score of 0.833. This superior performance
likely stems from \ac{YOLO}’s capability to extract fine-grained spatial features in high-resolution images, combined
with its efficient end-to-end object-centric architecture.

The \textbf{\ac{ResNet}} and \textbf{\ac{ViT}} models also demonstrated strong performance, with F1-scores of 0.810 and
0.790, respectively. \ac{ResNet} maintained a slightly higher accuracy (0.879) than \ac{ViT} (0.876), suggesting
that deep convolutional features still offer competitive advantages. Meanwhile, \ac{ViT}’s transformer-based design
effectively captured semantic dependencies, reinforcing its value for tasks involving subtle symbol variations.

In contrast, \textbf{EfficientNet} underperformed in this task, with an F1-score of 0.709 and a macro recall of
0.705, indicating difficulty in generalizing across all classes. The \textbf{Moondream} model also showed weak
results (F1-score of 0.355), despite a moderately high recall (0.601), suggesting overprediction of certain classes.

The \textbf{OpenCLIP zero-shot} model performed poorly overall, with an accuracy of 0.274 and F1-score of 0.265.
These results reinforce that zero-shot approaches struggle when applied to highly specialized tasks like Nazi symbol
recognition without task-specific adaptation or supervised training.

When OpenCLIP embeddings were combined with traditional classifiers, performance varied widely.
\textbf{OpenCLIP + KNeighborsClassifier} (F1-score: 0.775) and \textbf{OpenCLIP + \ac{SVC}} (F1-score: 0.784) were
the most effective, indicating that neighborhood-based and kernel-based methods are capable of partially exploiting
OpenCLIP’s semantic features. Similarly, \textbf{OpenCLIP + SGDClassifier} and
\textbf{OpenCLIP + GradientBoostingClassifier} achieved solid F1-scores of 0.718 and 0.703, respectively.

Conversely, ensemble-based models such as \textbf{RandomForest}, \textbf{Bagging}, and
\textbf{DecisionTreeClassifier} performed poorly, with F1-scores around or below 0.200. These results suggest that
these models may not align well with the high-dimensional structure of OpenCLIP embeddings, potentially due to
overfitting or limited ability to model fine-grained inter-class distinctions.

Overall, end-to-end models such as \ac{YOLO}, \ac{ResNet}, and \ac{ViT} clearly outperform both zero-shot approaches and
OpenCLIP embedding pipelines in the multi-class setting. These findings emphasize the importance of strong visual
feature extraction and task-aligned architectures when handling classification tasks involving subtle visual
differences across many categories.

\subsection{Comparison with Binary Classification Performance}

Comparing the results from the binary (~\autoref{tab:binary-results}) and multi-class
(~\autoref{tab:multi-class-results}) classification tasks reveals consistent trends in model effectiveness, but also
highlights some key differences driven by task complexity.

In both settings, \textbf{\ac{YOLO}} and \textbf{\ac{ViT}} emerged as top-performing \ac{DL} models. However, their
relative strengths became more pronounced in the binary task. For instance, \ac{ViT} achieved an F1-score of 0.901 and
\ac{YOLO} reached 0.857 in the binary case, whereas their multi-class F1-scores dropped to 0.790 and 0.833,
respectively. This drop reflects the increased difficulty of distinguishing between a larger number of visually
similar symbols.

The benefit of combining \textbf{OpenCLIP embeddings with classical classifiers} was more significant in the binary
setting. In that task, models like \textbf{OpenCLIP + \ac{SVC}} and \textbf{OpenCLIP + KNeighborsClassifier} achieved
F1-scores of 0.939 and 0.915 respectively, outperforming nearly all other approaches. However, in the multi-class
scenario, while still strong (F1-scores of 0.784 and 0.775 respectively), their advantage diminished, likely due to
increased inter-class ambiguity and the limitations of traditional classifiers in modeling fine-grained boundaries
across many classes.

Interestingly, some ensemble models (e.g., GradientBoosting, VotingClassifier) showed reasonable binary
classification performance (F1-scores around 0.840–0.910), but generally degraded in the multi-class case (F1-scores
below 0.710), possibly due to overfitting or insufficient class-separation power.

The \textbf{Moondream} and \textbf{OpenCLIP zero-shot} models performed poorly in both tasks, with Moondream showing
a pathological tendency to predict the positive class in binary classification (perfect recall but extremely low
precision), and generally low macro F1-scores in multi-class classification. These outcomes highlight the risks of
deploying zero-shot or generative models in specialized symbol recognition tasks without task-specific training.

Overall, binary classification appears more forgiving, allowing classical classifiers with pretrained embeddings to
compete closely with end-to-end \ac{DL} models. In contrast, the multi-class task places a heavier burden on
precise visual discrimination and inter-class separability, where deep models like \ac{YOLO} and \ac{ViT} retain a clear
edge. This underscores the importance of matching model complexity and architecture with task granularity and class
diversity.

\subsection{Multi-Class Classification Confusion Matrix Analysis}

The confusion matrices for the multi-class classification task reveal distinct patterns of model performance across
classes, reflecting varying degrees of class separation and misclassification.

\textbf{\ac{ResNet}} (see ~\autoref{fig:resnet-confusion-matrix-multiple}) shows strong true positive rates for
high-frequency classes such as \textit{Swastika} (808 correct predictions) and \textit{Siegrune} (247 correct), but
also notable confusion between visually similar classes. For example, \textit{SS Skull} is frequently misclassified as
\textit{Siegrune} (36 instances), while \textit{Siegrune} mislabels \textit{SS Skull} less often (5 instances). This
asymmetric confusion highlights the difficulty \ac{ResNet} faces in discriminating fine-grained symbol variations.
Low-frequency classes like \textit{Nazi Salute} perform poorly, with very few correct predictions.

\textbf{\ac{YOLO}} (see ~\autoref{fig:yolo-confusion-matrix-multiple}) generally maintains competitive accuracy for
dominant classes, correctly identifying 834 \textit{Swastika} instances and 222 \textit{Siegrune} samples. Confusion
between \textit{SS Skull} and \textit{Siegrune} remains substantial but somewhat more balanced compared to
\ac{ResNet}, with 25 \textit{SS Skull} instances misclassified as \textit{Siegrune} and 19 misclassifications in the
opposite direction. This indicates improved inter-class separation but ongoing challenges for these similar classes.

\textbf{EfficientNet} (see ~\autoref{fig:efficient-confusion-matrix-multiple}) exhibits a somewhat different confusion
profile, with lower misclassification of \textit{SS Skull} as \textit{Siegrune} (15 instances), but higher
misclassification of \textit{Siegrune} as \textit{SS Skull} (40 instances), showing an inverse asymmetry compared to
\ac{ResNet}. The true positive count for \textit{Siegrune} is comparatively lower (174), suggesting less confident
classification in this region of the feature space. Low-frequency classes continue to be problematic, with
\textit{Nazi Salute} achieving only a single correct prediction.

\textbf{\ac{ViT}} (see ~\autoref{fig:vit-confusion-matrix-multiple}) marginally improves performance for certain classes, such as
\textit{SS Skull} and \textit{Hitler}, with 188 and 175 correct predictions respectively. Confusion between
\textit{SS Skull} and \textit{Siegrune} remains, with 19 \textit{SS Skull} samples misclassified as \textit{Siegrune} and 32 in the
reverse direction, indicating persistent challenges in distinguishing subtle differences despite the transformer
architecture’s strengths.

\textbf{OpenCLIP Encoder combined with \ac{SVC}} (see ~\autoref{fig:openclip-encoder-confusion-matrix-multiple})
achieves a balanced confusion pattern. It correctly identifies 232 \textit{Siegrune} and 167 \textit{SS Skull}
instances. The confusion between \textit{SS Skull} and \textit{Siegrune} is distributed moderately (25 and 17
misclassifications respectively), showing a more symmetric error pattern compared to other models. This suggests
better inter-class separation in this challenging pair, although some misclassification persists. Additionally, the
model excels at maintaining tight clusters around true labels for high-frequency classes such as
\textit{Swastika} (823 correct) but struggles with low-frequency classes like \textit{Nazi Salute}, which has zero
correct predictions.

Overall, all models face difficulties with fine-grained classes like \textit{SS Skull} and \textit{Siegrune}, but
OpenCLIP+\ac{SVC} demonstrates a more balanced confusion profile. The results highlight the ongoing need to improve
discrimination for visually similar symbols, especially for classes with limited training data.

\section{Analysis of Underperforming Models}

While several models demonstrated strong classification performance, it is equally important to investigate those
that underperformed to understand their limitations and potential failure modes.

The \textbf{Moondream} model demonstrated systematic failure in the binary classification task, marked by an extreme
trade-off between recall (0.995) and precision (0.153), yielding an overall F1-score of only 0.265. This suggests it
overwhelmingly classified images as Nazi symbols, resulting in a high false positive rate. Unlike other models,
Moondream is primarily designed for vision-language tasks such as captioning and dialogue generation, not
fine-grained visual classification. Its generative nature and alignment safeguards may contribute to
overgeneralization or suppression of sensitive content, potentially causing it to indiscriminately signal the
presence of harmful symbols. In multi-class classification, its performance remained weak (F1-score: 0.355),
underscoring the model’s limited visual discrimination capability. These results highlight that even models with
strong multimodal language capabilities may perform poorly on classification tasks that demand high visual
specificity and consistent decision boundaries—particularly when not explicitly fine-tuned for such objectives.

The \textbf{OpenCLIP zero-shot} model underperformed across both binary and multi-class classification settings. In
the binary task, it achieved an F1-score of just 0.025, and in the multi-class setting, it reached only 0.265
F1-score and 0.274 accuracy. These results suggest that OpenCLIP’s generic semantic alignment—though effective for
open-ended retrieval or captioning—falls short for tasks requiring precise, domain-specific visual understanding.
The lack of supervised fine-tuning likely hindered the model’s ability to map subtle symbol variations to specific
class labels, especially in a setting involving ideologically charged and visually nuanced symbols. Moreover, its
pretrained text embeddings may not reflect the necessary granularity or cultural-historical specificity required to
distinguish, for example, a stylized swastika from unrelated geometric patterns. These outcomes reinforce the
critical need for task- and domain-specific adaptation when applying general-purpose foundation models to
specialized classification tasks.

Several \textbf{classical ensemble models} combined with OpenCLIP embeddings, such as the
\textbf{RandomForestClassifier}, exhibited counterintuitive behavior. While the binary classification results showed
extremely high precision (0.983), they suffered from near-total recall failure (0.097), resulting in a poor F1-score
of 0.177. This pattern suggests an overly cautious decision boundary, where the model rarely labeled images as
positive—possibly due to poor generalization in the high-dimensional OpenCLIP feature space. In the multi-class
setting, ensemble models such as Random Forest, Bagging, and Decision Tree failed to exceed 0.200 F1-score, likely
due to their relatively shallow decision structures and limited capacity to model the complex, fine-grained
relationships embedded in pretrained vision-language features. These findings indicate that while classical models
can perform well in low-dimensional or well-separated feature spaces, they struggle to generalize when paired with
dense, high-capacity representations unless further feature selection or dimensionality reduction is applied.

Overall, these underperformances highlight the importance of model-task alignment. Generative or zero-shot models
may offer generality, but without fine-tuning, they can fail catastrophically on specific classification tasks.
Likewise, conventional tree-based models may struggle to effectively utilize dense, semantic embeddings like those
from OpenCLIP, especially in tasks with fine-grained categories.

These failures provide a contrast to the following high-performing approaches, which better align with the
classification objectives.

\section{Analysis of High-Performing Models}

Across both binary and multi-class classification tasks, several models demonstrated exceptional performance,
providing insight into the characteristics of effective architectures for Nazi symbol recognition.

The highest-performing binary classifier was the combination of
\textbf{OpenCLIP embeddings with a \ac{SVC}}, which achieved an F1-score of 0.939, paired
with a high precision (0.948) and recall (0.929). This suggests that the kernel-based decision boundaries of \ac{SVC} are
well-suited to partitioning OpenCLIP’s semantically rich, high-dimensional embedding space. Its confusion matrix
reveals both low false positive and false negative rates, indicating stable and balanced performance even under
class imbalance. Similarly, the \textbf{OpenCLIP + \ac{LR}} model achieved the best \ac{AUC} (1.000) and \ac{AP} (
0.964), highlighting its strength in probabilistic calibration. The fact that even a linear model performs so well
when paired with OpenCLIP embeddings underlines the quality of these pretrained representations, which already
encode strong semantic separability between Nazi and non-Nazi symbols. Together, these results demonstrate that
lightweight classifiers—when paired with powerful pretrained vision-language embeddings—can deliver highly
competitive performance on specialized tasks, without requiring deep end-to-end training.

In the multi-class classification setting, \textbf{\ac{YOLO}} emerged as the top-performing model, achieving a macro
F1-score of 0.833. Unlike purely classification-based architectures, \ac{YOLO}’s detection-oriented structure inherently
emphasizes spatial localization and region-level features, which may be especially advantageous for distinguishing
fine-grained visual details among similar-looking symbol classes. Its strong per-class performance on high-frequency
symbols like \textit{Swastika} and \textit{Siegrune} suggests it effectively captures both global and local
patterns. Moreover, \ac{YOLO}’s anchor-based detection approach may offer a form of implicit segmentation, aiding the
model in isolating relevant symbol regions from noisy or cluttered image backgrounds—a recurring challenge in
real-world deployment contexts.

Both \textbf{ResNet} and \textbf{\ac{ViT}} demonstrated consistently strong performance across binary and multi-class
tasks, with F1-scores above 0.79 in all settings. ResNet’s deep convolutional layers facilitated robust extraction
of low- to mid-level visual features, while its residual connections likely improved gradient flow and optimization
stability during training. \ac{ViT}, in contrast, relies on a patch-based self-attention mechanism, which enables it to
model long-range dependencies and complex compositional structures—a key advantage when symbols contain subtle
stylistic variations. Their strong multi-task performance reinforces the effectiveness of end-to-end deep learning
architectures when trained on sufficient, well-labeled domain data. These models serve as reliable baselines for
further exploration of task-specific adaptations or multimodal extensions.

Among embedding-based approaches, \textbf{OpenCLIP + KNeighborsClassifier} also stood out with strong binary (
F1-score: 0.915) and competitive multi-class (F1-score: 0.775) performance. The effectiveness of KNN suggests that
OpenCLIP embeddings preserve semantic distances in a way that aligns well with class membership—enabling
nearest-neighbor comparisons to retrieve visually and semantically similar symbols. This model’s strength may also
stem from its non-parametric nature, which avoids assumptions about the shape of decision boundaries and instead
relies on local geometric structure in the embedding space. Such behavior could be particularly valuable in
classification settings where labeled examples are sparse but clustered, and where explainability through exemplar
retrieval is desirable.

Across high-performing models, a common theme emerges: the combination of rich visual representations—whether
through end-to-end convolutional or transformer architectures, or pretrained embeddings like OpenCLIP—with
well-matched classification strategies, yields robust symbol classification. These results emphasize the value of
pairing strong feature extractors with appropriately complex decision functions, tailored to the geometry of the
input space and the semantic structure of the task.

Beyond individual model metrics, broader themes emerge across tasks and models, offering insight into effective
strategies for symbol classification.

\section{Key Insights and Thematic Patterns}
\label{subsec:key-insights}

Analysis across both binary and multi-class tasks revealed several recurring patterns that shed light on the design
trade-offs and performance dynamics of different model classes.

\paragraph{Hybrid architectures outperform raw foundation models.}
One of the most consistent findings was the superior performance of OpenCLIP embeddings when combined with classical
machine learning classifiers such as \ac{SVC}, \ac{LR}, or \ac{KNN}. These hybrid models outperformed both
traditional \ac{CNN}s and foundation models used in isolation (e.g., zero-shot OpenCLIP or Moondream). The results
suggest that while pretrained vision-language embeddings capture rich semantic structure, they require tailored
downstream classifiers to translate these representations into accurate symbol predictions. This modular approach
also offers flexibility in tuning precision-recall trade-offs and model interpretability.

\paragraph{Transformer-based models excel in nuanced visual tasks.}
Vision Transformer (\ac{ViT}) models demonstrated consistently strong results, benefiting from global
self-attention mechanisms that allow for nuanced discrimination between visually similar symbols. Compared to \ac{CNN}s
like ResNet or EfficientNet, \ac{ViT} showed greater robustness to stylization, partial occlusion, and background noise.
This suggests that transformer architectures may be better suited for iconographic classification tasks that require
compositional understanding and spatial context awareness.

\paragraph{Model architecture impacts suitability for deployment.}
The classification-specific variant of the YOLO model—developed by the YOLO team—achieved the best performance in
the multi-class classification task, with the highest accuracy (0.884), precision (0.865), recall (0.818), and
F1-score (0.833). These results highlight the strength of streamlined, efficient architectures in real-world
settings where both performance and inference speed matter. Unlike traditional object detectors, this version of
YOLO is natively optimized for classification tasks. Its ability to maintain \ac{SotA} accuracy while
remaining computationally lightweight makes it especially well-suited for deployment in constrained environments or
applications requiring real-time symbol recognition.

\paragraph{Generalization to historical and stylized symbols remains a challenge.}
Several models struggled when presented with visually degraded, stylized, or low-contrast images typical of
historical archives or user-generated content. Despite strong performance on curated datasets, generalization gaps
remain, particularly for rare or obscure Nazi symbols. This points to the need for domain adaptation strategies such
as contrastive pretraining, synthetic augmentation, or inclusion of diverse training sources to enhance model
robustness.

\paragraph{Symbol frequency imbalance affects recall and stability.}
As in many multi-class classification settings, class imbalance significantly impacted model performance. Rare
symbol classes, such as the \textit{Black Sun}, exhibited low recall and unstable predictions,
even in otherwise high-performing models. Addressing this will likely require targeted techniques such as few-shot
learning, oversampling, or curriculum-based training to ensure equitable performance across all classes.

Together, these insights highlight the trade-offs between model complexity, interpretability, and task
suitability—informing how future systems can balance accuracy with practicality and ethical deployment needs.

\section{Implications for Real-World Deployment}
\label{subsec:deployment-implications}

The findings suggest that hybrid approaches—combining deep feature extraction with classical \ac{ML}
classifiers—offer the most promise for deployment in real-world systems tasked with binary classification of Nazi
symbols in images. In particular, OpenCLIP combined with \ac{SVC} consistently achieved the highest F1-score in the
binary classification task, indicating strong performance in balancing precision and recall when identifying rare and
sensitive symbols. These hybrid models also offer a favorable trade-off between accuracy, interpretability, and
computational efficiency—key considerations for safety-critical applications such as content moderation or archival
screening.

However, practical deployment requires addressing several challenges. These include handling ambiguous or
degraded images, unseen symbol variants, and the high cost of false positives in legal, historical, or moderation
contexts. The results suggest that a human-in-the-loop approach may be appropriate, where automated classification is
supplemented by expert review for uncertain or high-stakes cases.

Moreover, compliance and interpretability are critical. Simpler classifiers layered on OpenCLIP embeddings offer
greater transparency than end-to-end deep networks, which can support auditability, regulatory acceptance, and user
trust.

These deployment-oriented considerations directly support the research objective of developing practical, accurate, and
ethically responsible classification systems for real-world symbol recognition tasks.

\section{Limitations and Risks}
\label{subsec:limitations}

Despite promising results, several limitations and risks remain. Most notably, model performance on minority classes
was inconsistent, with a tendency to overfit to frequent, visually dominant symbols. This raises concerns around
fairness, completeness, and the system's ability to generalize to underrepresented symbol types in deployment.

Dataset bias also presents a notable risk. Much of the training data was synthetically curated or drawn from known
repositories, potentially omitting rare, obscure, or modern reinterpretations of Nazi iconography. As a result, the
system may misclassify or entirely miss symbols that deviate from the curated training distribution.

Another key limitation is the lack of contextual awareness. The models operate solely on visual information and do not
distinguish between educational, artistic, or propagandistic usage. While context modeling was beyond the scope of
this project, it is essential for responsible deployment. Without it, there is a risk of misinterpretation, overreach,
or harm—particularly in sensitive social, legal, or historical settings.

These limitations highlight the need for further work in data diversification, contextual modeling, and ethical
evaluation before deploying such systems at scale.
